\documentclass[12pt,a4paper]{report}

% ============================================================
% PACKAGES
% ============================================================
\usepackage[utf8]{inputenc}
\usepackage[T1]{fontenc}
\usepackage[french]{babel}
\usepackage{geometry}
\usepackage{graphicx}
\usepackage{hyperref}
\usepackage{fancyhdr}
\usepackage{titlesec}
\usepackage{array}
\usepackage{booktabs}
\usepackage{xcolor}
\usepackage{float}
\usepackage{enumitem}
\usepackage{setspace}
\usepackage{caption}
\usepackage{parskip}
\usepackage{tocloft}
\usepackage{emptypage}
\usepackage{microtype}

\geometry{
    top=2.5cm,
    bottom=2.5cm,
    left=3cm,
    right=2.5cm
}

% ============================================================
% COULEURS
% ============================================================
\definecolor{ensias-red}{RGB}{180,0,0}
\definecolor{ensias-dark}{RGB}{30,30,30}
\definecolor{gray-light}{RGB}{245,245,245}

% ============================================================
% STYLE DES CHAPITRES
% ============================================================
\titleformat{\chapter}[display]
  {\normalfont\huge\bfseries\color{ensias-dark}}
  {\chaptertitlename\ \thechapter}
  {20pt}
  {\Huge\bfseries}
\titlespacing*{\chapter}{0pt}{-10pt}{30pt}

% ============================================================
% EN-TÊTES ET PIEDS DE PAGE
% ============================================================
\pagestyle{fancy}
\fancyhf{}
\renewcommand{\headrulewidth}{0.4pt}
\fancyhead[L]{\small\nouppercase{\leftmark}}
\fancyhead[R]{\small ENSIAS}
\fancyfoot[L]{\small ENSIAS}
\fancyfoot[C]{\thepage}
\fancyfoot[R]{\small 2024-2025}

% ============================================================
% HYPERLIENS
% ============================================================
\hypersetup{
    colorlinks  = true,
    linkcolor   = ensias-dark,
    urlcolor    = ensias-red,
    citecolor   = ensias-dark,
    pdftitle    = {Système de Gestion des Dossiers Médicaux},
    pdfauthor   = {Fatima-Ezzahra Qasmi, Mohamed Khalil Ibouchna},
}

% ============================================================
% TABLE DES MATIÈRES -- style
% ============================================================
\renewcommand{\cftchapfont}{\bfseries}
\renewcommand{\cftsecfont}{\normalfont}
\renewcommand{\cftsubsecfont}{\normalfont}
\setlength{\cftbeforechapskip}{6pt}

% ============================================================
% DÉBUT DU DOCUMENT
% ============================================================
\begin{document}

% -----------------------------------------------------------
% PAGE DE GARDE
% -----------------------------------------------------------
\begin{titlepage}
\centering

\begin{figure}[h]
    \centering
    \begin{minipage}{0.35\textwidth}
        \centering
        \includegraphics[width=3.5cm]{Ensias logo.jpg}
    \end{minipage}
    \hfill
    \begin{minipage}{0.35\textwidth}
        \centering
        \includegraphics[width=3.5cm]{logo_um5.png}
    \end{minipage}
\end{figure}

\vspace{0.6cm}

{\scshape\Large École Nationale Supérieure}\\[0.2cm]
{\scshape\Large d'Informatique et d'Analyse des Systèmes}\\[0.2cm]
{\scshape\Large -- Rabat}

\vspace{0.8cm}
\rule{\textwidth}{1pt}\\[0.3cm]
{\large Projet de Fin d'Année}\\[0.2cm]
\rule{\textwidth}{1pt}

\vspace{1cm}

{\Huge\bfseries Système de Gestion des Dossiers}\\[0.4cm]
{\Huge\bfseries Médicaux :}\\[0.4cm]
{\Huge\bfseries Application des modèles d'IA}

\vspace{1.2cm}
\rule{\textwidth}{0.5pt}

\vfill

\begin{minipage}[t]{0.45\textwidth}
    \raggedright
    \textbf{Élèves :}\\[4pt]
    Fatima-Ezzahra \textsc{Qasmi}\\
    Mohamed Khalil \textsc{Ibouchna}
\end{minipage}
\hfill
\begin{minipage}[t]{0.45\textwidth}
    \centering
    \textbf{Encadrant :}\\[4pt]
    Pr.~Youness \textsc{Tabii}\\[0.6cm]
    \textbf{Jury :}\\[4pt]
    Prénom \textsc{Nom}\\
    Prénom \textsc{Nom}
\end{minipage}

\vspace{1cm}
{\large Année Universitaire 2024-2025}
\end{titlepage}

% -----------------------------------------------------------
% REMERCIEMENTS
% -----------------------------------------------------------
\chapter*{Remerciements}
\addcontentsline{toc}{chapter}{Remerciements}
\thispagestyle{fancy}

Ce projet de fin d'année représente bien plus qu'un simple exercice académique : il
incarne l'aboutissement d'une année riche en apprentissages, en découvertes et en défis.

Nous souhaitons exprimer notre sincère gratitude à notre encadrant \textbf{Pr.~Youness
TABII}, pour sa précieuse disponibilité, ses conseils avisés, et son accompagnement constant
tout au long de ce projet. Sa rigueur scientifique et sa vision claire ont été des sources
d'inspiration et de motivation pour mener à bien cette réalisation.

Nos remerciements s'adressent également à l'ensemble du corps professoral de
l'ENSIAS pour la qualité de l'enseignement et le cadre stimulant qu'ils nous ont offert
durant notre parcours.

Enfin, nous remercions nos familles et nos amis pour leur soutien moral inconditionnel
tout au long de cette aventure académique.

% -----------------------------------------------------------
% RÉSUMÉ
% -----------------------------------------------------------
\chapter*{Résumé}
\addcontentsline{toc}{chapter}{Résumé}
\thispagestyle{fancy}

Ce projet de fin d'année consiste en la conception et le développement d'un système web de gestion des dossiers médicaux destiné à une clinique, intégrant un module d'intelligence artificielle pour l'aide au diagnostic.\\

L'application centralise la gestion des patients, des rendez-vous, des consultations et des prescriptions.\\

Lors d'une consultation, le médecin sélectionne les symptômes observés chez le patient, et le système retourne automatiquement les maladies les plus probables accompagnées de leur score de probabilité, grâce à un modèle de classification entraîné sur le dataset \textbf{Disease Prediction Using Machine Learning} (132 symptômes, 42 maladies).\\

L'architecture repose sur une approche microservices : un backend Spring Boot (Java) connecté à une base PostgreSQL pour la logique métier, un microservice Python (FastAPI) dédié au modèle d'IA, et une interface React pour les utilisateurs. La conception suit les standards UML (cas d'utilisation, classes, séquences).


\medskip
\noindent\textbf{Mots-clés :} Gestion des dossiers médicaux, intelligence artificielle,
prédiction de maladies, Spring Boot, React, FastAPI, microservices, JWT.

% -----------------------------------------------------------
% LISTE DES ABRÉVIATIONS
% -----------------------------------------------------------
\chapter*{Liste des Abréviations}
\addcontentsline{toc}{chapter}{Liste des Abréviations}
\thispagestyle{fancy}

\begin{tabular}{p{2.5cm}p{11cm}}
\textbf{API}     & Application Programming Interface \\[4pt]
\textbf{CORS}    & Cross-Origin Resource Sharing \\[4pt]
\textbf{CRUD}    & Create, Read, Update, Delete \\[4pt]
\textbf{DTO}     & Data Transfer Object \\[4pt]
\textbf{ENSIAS}  & École Nationale Supérieure d'Informatique et d'Analyse des Systèmes \\[4pt]
\textbf{IA}      & Intelligence Artificielle \\[4pt]
\textbf{JPA}     & Jakarta Persistence API \\[4pt]
\textbf{JWT}     & JSON Web Token \\[4pt]
\textbf{ML}      & Machine Learning \\[4pt]
\textbf{ORM}     & Object-Relational Mapping \\[4pt]
\textbf{PDF}     & Portable Document Format \\[4pt]
\textbf{RBAC}    & Role-Based Access Control \\[4pt]
\textbf{REST}    & REpresentational State Transfer \\[4pt]
\textbf{SQL}     & Structured Query Language \\[4pt]
\textbf{UML}     & Unified Modeling Language \\[4pt]
\end{tabular}

% -----------------------------------------------------------
% TABLE DES MATIÈRES
% -----------------------------------------------------------
\clearpage
\tableofcontents
\clearpage

% -----------------------------------------------------------
% TABLE DES FIGURES
% -----------------------------------------------------------
\listoffigures
\addcontentsline{toc}{chapter}{Table des figures}
\clearpage

% -----------------------------------------------------------
% LISTE DES TABLEAUX
% -----------------------------------------------------------
\listoftables
\addcontentsline{toc}{chapter}{Liste des tableaux}
\clearpage

% -----------------------------------------------------------
% INTRODUCTION GÉNÉRALE
% -----------------------------------------------------------
\chapter*{Introduction générale}
\addcontentsline{toc}{chapter}{Introduction générale}
\thispagestyle{fancy}

La transformation numérique du secteur de la santé s'est imposée comme un levier majeur d'amélioration de la qualité des soins. Le passage du dossier médical papier vers le dossier médical électronique centralise l'information clinique, réduit les erreurs administratives et fournit une base exploitable pour des analyses prédictives. Parallèlement, l'émergence de l'intelligence artificielle ouvre de nouvelles perspectives, notamment en matière d'aide à la décision diagnostique à partir des symptômes observés.\\

Cependant, les cliniques de petite et moyenne taille rencontrent souvent des difficultés à se doter de solutions logicielles à la fois simples, fiables et intégrant ces technologies. Les outils existants sont fréquemment coûteux ou inadaptés à une structure mono-site, ce qui maintient des pratiques manuelles limitant l'efficacité opérationnelle et la qualité du suivi patient.\\

C'est dans cette problématique que s'inscrit le présent projet de fin d'année, intitulé « Système de gestion des dossiers médicaux : application des modèles d'intelligence artificielle », réalisé dans le cadre de la première année du cycle d'ingénieur en Génie Logiciel à l'ENSIAS, sous la supervision du \textbf{Pr.Youness Tabii}. L'objectif est de concevoir et développer un site web destinée à une clinique unique, couvrant l'ensemble du cycle de prise en charge d'un patient et intégrant un module d'aide au diagnostic basé sur un modèle d'apprentissage automatique.\\

Le système distingue trois catégories d'utilisateurs (médecin, secrétaire, patient), chacune avec des fonctionnalités adaptées à son rôle. Sa fonctionnalité différenciante réside dans le module d'intelligence artificielle : lors d'une consultation, le médecin sélectionne les symptômes observés chez le patient parmi une liste de cent trente-deux symptômes, et le système retourne les pathologies les plus probables accompagnées de leur score de probabilité, parmi quarante-deux maladies couvertes par le modèle. Cet outil ne se substitue pas au jugement clinique, mais enrichit la réflexion diagnostique du praticien.\\

Le développement de l’application a été réalisé à l’aide des technologies Spring Boot (Java) en backend, une base de données PostgreSQL, une interface utilisateur en React, et un microservice d'IA développé en Python avec FastAPI. La conception suit les standards UML (cas d'utilisation, classes, séquences) , permettant de valider chaque brique fonctionnelle avant d'enrichir progressivement le système.\\

Ce rapport présente les différentes étapes du projet, de l’analyse du besoin à la mise
en œuvre de la solution, en passant par la conception technique, la modélisation, le développement et les tests fonctionnels.


% ===========================================================
% CHAPITRE 1 -- CONTEXTE GÉNÉRAL
% ===========================================================
\chapter{Contexte Général}

\section*{Introduction}
\addcontentsline{toc}{section}{Introduction}

Dans le secteur de la santé, la gestion manuelle des dossiers médicaux représente un
frein majeur à l'efficacité clinique et à la qualité du suivi patient. La multiplicité des
intervenants, la sensibilité des données traitées et les exigences croissantes en matière de
traçabilité rendent indispensable l'adoption de solutions numériques adaptées. Ce chapitre
présente le contexte dans lequel s'inscrit notre projet, la problématique à résoudre, les
objectifs visés et la méthodologie de développement adoptée.

% ---
\section{Problématique}

Les cliniques de petite et moyenne taille font face à plusieurs défis dans la gestion
quotidienne de leurs opérations :

\begin{itemize}
    \item \textbf{Absence de centralisation des dossiers médicaux :} les informations
    relatives aux patients (consultations, ordonnances, analyses, historique) sont
    dispersées sur support papier ou dans des fichiers non structurés, rendant
    difficile toute consultation rapide.

    \item \textbf{Processus de prise de rendez-vous inefficace :} la gestion manuelle
    des plannings engendre des erreurs, des doublons et une mauvaise utilisation du
    temps médical.

    \item \textbf{Absence d'aide au diagnostic :} les médecins ne disposent d'aucun
    outil d'aide à la décision exploitant l'ensemble des symptômes observés pour
    orienter vers les pathologies les plus probables.

    \item \textbf{Manque d'accès autonome pour les patients :} les patients ne
    disposent pas d'un espace personnel pour consulter leurs dossiers, résultats ou
    historiques de consultations.

    \item \textbf{Coût prohibitif des solutions du marché :} les logiciels médicaux
    disponibles sont souvent conçus pour de grandes structures hospitalières et
    représentent un investissement inaccessible pour une clinique mono-site.
\end{itemize}

Face à ces constats, la question centrale se pose : \textit{Comment concevoir une
application web complète, sécurisée et accessible, permettant à la Clinique Union de gérer
efficacement ses dossiers médicaux tout en intégrant un module d'aide au diagnostic par
intelligence artificielle~?}

% ---
\section{Objectifs}

\subsection{Objectifs principaux}

\begin{itemize}
    \item \textbf{Centraliser les dossiers médicaux :} remplacer les supports papier
    par une plateforme numérique centralisée, accessible et sécurisée.

    \item \textbf{Automatiser la gestion des rendez-vous :} permettre aux patients
    de prendre leurs rendez-vous en ligne et au secrétariat de gérer les plannings.

    \item \textbf{Intégrer un module d'aide au diagnostic par IA :} lors de chaque
    consultation, le médecin obtient une prédiction de maladie avec son score de
    probabilité à partir des symptômes sélectionnés.

    \item \textbf{Offrir un assistant conversationnel médical :} mettre à disposition
    des patients un chatbot propulsé par Gemini pour répondre à leurs questions de
    santé et les orienter.

    \item \textbf{Assurer la sécurité des données :} protéger les données sensibles
    des patients via une authentification basée sur des tokens JWT et une gestion
    fine des rôles.
\end{itemize}

\subsection{Fonctionnalités clés de la solution}

\begin{itemize}
    \item Inscription et authentification sécurisée par rôle (Médecin, Patient,
    Secrétaire) avec tokens JWT.

    \item Gestion complète des dossiers médicaux : consultations, ordonnances,
    documents PDF.

    \item Prise et gestion des rendez-vous en ligne avec validation par le
    secrétariat.

    \item Sélection interactive de symptômes et prédiction de maladie en temps
    réel via le modèle Random Forest (132 symptômes, 41 maladies).

    \item Assistant chatbot médical (Gemini) disponible 24h/24 pour les patients.

    \item Téléchargement et gestion des documents médicaux (ordonnances, analyses
    PDF).

    \item Tableaux de bord adaptés à chaque profil utilisateur.
\end{itemize}

\subsection{Innovation et différenciation}

\begin{itemize}
    \item \textbf{Approche microservices :} découplage entre la logique métier
    (Spring Boot), le modèle d'IA (FastAPI/Python) et l'interface utilisateur
    (React).

    \item \textbf{Modèle ML embarqué :} un RandomForestClassifier (accuracy 99,9\,\%)
    entraîné sur le dataset \textit{Disease Prediction Using Machine Learning} est
    exposé comme microservice REST.

    \item \textbf{Chatbot médical Gemini :} intégration d'un modèle de langage de
    grande taille pour offrir une assistance médicale conversationnelle
    personnalisée.

    \item \textbf{Architecture sécurisée :} authentification stateless par JWT,
    contrôle d'accès basé sur les rôles (RBAC), hachage BCrypt des mots de passe.
\end{itemize}

% ---
\section{Méthodologie et approche projet}

\subsection{Démarche de développement}

Le projet adopte une \textbf{méthodologie agile itérative}, organisée en phases
successives permettant une livraison progressive des fonctionnalités et une intégration
continue entre les composants.

Les grandes étapes sont les suivantes :
\begin{enumerate}
    \item \textbf{Analyse des besoins :} identification des acteurs, des cas
    d'utilisation et des exigences fonctionnelles et non fonctionnelles.

    \item \textbf{Conception :} modélisation UML (classes, séquences, cas
    d'utilisation) et conception de l'architecture microservices.

    \item \textbf{Développement itératif :} implémentation par couches -- backend,
    microservice IA, frontend -- avec tests d'intégration à chaque étape.

    \item \textbf{Tests fonctionnels :} validation de chaque fonctionnalité sur
    l'ensemble des profils utilisateurs.

    \item \textbf{Livraison et documentation :} rédaction du rapport et mise en
    production locale de l'ensemble des services.
\end{enumerate}

\section*{Conclusion}
\addcontentsline{toc}{section}{Conclusion}

Ce chapitre a présenté le contexte général du projet, la problématique qui lui donne
naissance, les objectifs visés et la méthodologie adoptée. La nécessité d'une solution
numérique intégrant l'intelligence artificielle pour la gestion des dossiers médicaux est
clairement établie. Le chapitre suivant s'intéressera à l'analyse détaillée des besoins et
à la conception de la solution.

% ===========================================================
% CHAPITRE 2 -- ANALYSE DES BESOINS ET CONCEPTION
% ===========================================================
\chapter{Analyse des besoins et Conception}

\section*{Introduction}
\addcontentsline{toc}{section}{Introduction}

L'objectif de ce chapitre est de définir précisément l'ensemble des exigences auxquelles
devra répondre notre application web. À travers une analyse fonctionnelle et
organisationnelle, nous identifions les acteurs principaux, les besoins spécifiques, ainsi
que les contraintes techniques. Cette étude est formalisée par des diagrammes UML couvrant
les cas d'utilisation, les classes métier et les séquences d'interaction.

% ---
\section{Liste des acteurs}

Le système distingue trois types d'utilisateurs humains, chacun disposant d'un espace
et de droits propres à son rôle :

\begin{description}[leftmargin=2cm, style=nextline]
    \item[\textbf{Médecin}]
    Il consulte et crée des dossiers médicaux, enregistre des consultations, sélectionne
    des symptômes pour obtenir une prédiction IA, téléverse des documents (ordonnances,
    analyses PDF) et gère son planning de rendez-vous.

    \item[\textbf{Patient}]
    Il consulte son dossier médical, ses consultations passées, ses documents médicaux,
    prend des rendez-vous en ligne et interagit avec l'assistant chatbot IA.

    \item[\textbf{Secrétaire}]
    Il gère le planning des rendez-vous, valide ou refuse les demandes des patients, et
    assiste les médecins dans la gestion administrative de la clinique.
\end{description}

% ---
\section{Besoins fonctionnels}

Voici les principaux besoins fonctionnels identifiés pour notre application :

\paragraph{Gestion de l'authentification}
\begin{itemize}
    \item Inscription des utilisateurs avec validation des données (email unique,
    mot de passe haché avec BCrypt).
    \item Connexion sécurisée par email/mot de passe avec génération d'un token JWT
    (durée de validité : 24~heures).
    \item Contrôle d'accès basé sur le rôle : chaque rôle n'accède qu'à son espace
    dédié.
\end{itemize}

\paragraph{Gestion des rendez-vous}
\begin{itemize}
    \item Prise de rendez-vous en ligne par le patient avec sélection du médecin,
    de la date et du motif.
    \item Validation, modification ou annulation par la secrétaire.
    \item Consultation du planning par le médecin (statuts : PLANIFIÉ, CONFIRMÉ,
    ANNULÉ).
\end{itemize}

\paragraph{Gestion des dossiers médicaux}
\begin{itemize}
    \item Création automatique d'un dossier médical lors de la première consultation.
    \item Enregistrement des consultations : motif, diagnostic, ordonnance, symptômes.
    \item Téléversement et téléchargement de documents PDF (ordonnances, analyses).
\end{itemize}

\paragraph{Module d'aide au diagnostic IA}
\begin{itemize}
    \item Sélection interactive des symptômes (autocomplete sur 132 symptômes).
    \item Appel automatique au microservice FastAPI lors de la création d'une
    consultation.
    \item Affichage de la maladie prédite, du score de confiance et du top-3 des
    pathologies probables.
\end{itemize}

\paragraph{Assistant chatbot médical}
\begin{itemize}
    \item Interface conversationnelle pour les patients.
    \item Réponses générées par le modèle Gemini (via OpenRouter).
    \item Historique de conversation maintenu pendant la session.
\end{itemize}

% ---
\section{Besoins non fonctionnels}

\begin{itemize}
    \item \textbf{Sécurité :} authentification stateless par JWT, mots de passe
    hachés avec BCrypt, CORS configuré pour les origines autorisées.

    \item \textbf{Performance :} temps de réponse de l'API $<$ 500\,ms pour les
    opérations CRUD, $<$ 3\,s pour la prédiction IA.

    \item \textbf{Disponibilité :} les trois microservices (backend, IA, frontend)
    peuvent être démarrés et arrêtés indépendamment.

    \item \textbf{Scalabilité :} architecture microservices permettant de scaler le
    service IA indépendamment du backend métier.

    \item \textbf{Maintenabilité :} code Java structuré en couches (Controller,
    Service, Repository), annotations Lombok pour réduire le code boilerplate.

    \item \textbf{Compatibilité :} interface responsive, compatible avec les
    navigateurs modernes (Chrome, Firefox, Edge).
\end{itemize}

% ---
\section{Choix de la méthode UML}

La modélisation de notre système suit les standards \textbf{UML~2.x} pour les raisons
suivantes :

\begin{itemize}
    \item \textbf{Universalité :} UML est la norme internationale de modélisation
    des systèmes orientés objet, alignée avec notre architecture Java/Spring Boot.

    \item \textbf{Expressivité :} les différents types de diagrammes couvrent toutes
    les perspectives du système (fonctionnelle, structurelle, comportementale).

    \item \textbf{Alignement avec JPA/Hibernate :} le diagramme de classes se
    traduit directement en entités JPA et en tables PostgreSQL via la stratégie
    d'héritage \textbf{JOINED}.
\end{itemize}

Nous avons utilisé les diagrammes suivants : diagramme de cas d'utilisation (vue
fonctionnelle globale), diagramme de classes (modèle structurel) et diagrammes de
séquences (flux d'interactions pour les scénarios clés).

% ---
\section{Diagramme de cas d'utilisation}

\begin{figure}[H]
    \centering
    \includegraphics[width=0.82\textwidth]{duc.png}
    \caption{Diagramme de cas d'utilisation global du système}
    \label{fig:usecase}
\end{figure}

Le diagramme présente les interactions entre les trois acteurs et les
fonctionnalités du système. Il illustre que :
\begin{itemize}
    \item \textbf{Le Patient} peut s'authentifier, consulter son dossier médical,
    prendre un rendez-vous et interagir avec le chatbot IA.

    \item \textbf{Le Médecin} peut s'authentifier, créer une consultation avec
    sélection de symptômes, obtenir une prédiction IA et gérer les documents.

    \item \textbf{La Secrétaire} peut s'authentifier, gérer les rendez-vous et
    enregistrer de nouveaux patients.
\end{itemize}

% ---
\section{Diagramme de classes}

Le diagramme de classes modélise les entités métier du système. L'héritage entre
\texttt{Utilisateur} et ses sous-classes est implémenté en JPA via la stratégie
\textbf{JOINED} (une table par classe avec jointure sur l'identifiant).

\begin{figure}[H]
    \centering
    \makebox[\textwidth][c]{\includegraphics[width=1.18\textwidth]{diagramme_classes.png}}
    \caption{Diagramme de classes du système}
    \label{fig:classes}
\end{figure}

Les principales entités sont les suivantes :

\begin{itemize}
    \item \textbf{Utilisateur} (abstraite) : id, nom, prénom, email, motDePasse,
    role, dateCreation.
    \item \textbf{Patient} (hérite d'Utilisateur) : cin, dateNaissance, téléphone.
    \item \textbf{Médecin} (hérite d'Utilisateur) : spécialité.
    \item \textbf{Secrétaire} (hérite d'Utilisateur) : numéro de bureau.
    \item \textbf{RendezVous} : dateHeure, motif, statut, patient, médecin.
    \item \textbf{DossierMedical} : patient, liste de consultations.
    \item \textbf{Consultation} : dateConsultation, motif, diagnostic, ordonnanceNotes,
    symptômes, prédictionIA.
    \item \textbf{PredictionIA} : maladiePrédite, probabilité, symptômesAnalysés.
    \item \textbf{Document} : nomFichier, contentType, taille, type (ORDONNANCE /
    ANALYSE).
\end{itemize}

% ---
\section{Diagrammes de séquences}

\subsection*{Scénario 1 -- Authentification JWT}

Lors de la connexion, le frontend envoie les identifiants au backend. Spring Security
vérifie les credentials via l'AuthenticationManager, charge l'utilisateur depuis PostgreSQL
et, en cas de succès, génère un token JWT retourné au client.

\begin{figure}[H]
    \centering
    \includegraphics[width=0.9\textwidth]{seq_authentification.png}
    \caption{Diagramme de séquence -- Authentification JWT}
    \label{fig:seq_login}
\end{figure}

\subsection*{Scénario 2 -- Création d'une consultation avec prédiction IA}

Lors de la création d'une consultation, le backend Spring Boot appelle automatiquement
le microservice FastAPI avec la liste des symptômes. Le modèle RandomForest retourne la
maladie prédite et le score de confiance, qui sont persistés dans la table
\texttt{predictions\_ia}.

\begin{figure}[H]
    \centering
    \includegraphics[width=\textwidth]{seq_consultation_ia.png}
    \caption{Diagramme de séquence -- Consultation avec prédiction IA}
    \label{fig:seq_consultation}
\end{figure}

\subsection*{Scénario 3 -- Gestion des rendez-vous}

Le patient soumet une demande de rendez-vous. Le service vérifie la disponibilité du
médecin : si le créneau est libre, le rendez-vous est enregistré ; sinon, une exception
est levée (\texttt{409 Conflict}). La secrétaire valide ensuite depuis son tableau de bord.

\begin{figure}[H]
    \centering
    \includegraphics[width=\textwidth]{seq_rendezvous.png}
    \caption{Diagramme de séquence -- Gestion des rendez-vous}
    \label{fig:seq_rdv}
\end{figure}

\section*{Conclusion}
\addcontentsline{toc}{section}{Conclusion}

Ce chapitre a formalisé les besoins fonctionnels et non fonctionnels du système et
fourni une modélisation UML complète. Le diagramme de classes constitue le socle du
schéma de base de données généré automatiquement par Hibernate. Le chapitre suivant
détaille la mise en \oe{}uvre technique de ces spécifications.

% ===========================================================
% CHAPITRE 3 -- RÉALISATION
% ===========================================================
\chapter{Réalisation}

\section*{Introduction}
\addcontentsline{toc}{section}{Introduction}

Après avoir bien défini les besoins et la conception du système, ce chapitre décrit la
mise en \oe{}uvre technique. Nous présentons d'abord les outils et technologies retenus,
puis détaillons l'implémentation des fonctionnalités principales au travers des interfaces
utilisateur de chaque profil.

% ---
\section{Outils et Technologies utilisés \& Méthodologie}

\subsection{Technologies}

\subsubsection{Spring Boot -- Backend}

\begin{figure}[H]
    \centering
    \includegraphics[width=5cm]{logo_spring_boot.png}
    \caption{Spring Boot}
    \label{fig:springboot}
\end{figure}

\textbf{Spring Boot 3.4.5} (Java 21) constitue le c\oe{}ur du backend. Il expose une
API REST sécurisée, gère la persistance des données via JPA/Hibernate et assure
l'authentification par JWT. Les modules Spring utilisés sont :

\begin{itemize}
    \item \texttt{spring-boot-starter-web} : exposition des endpoints REST.
    \item \texttt{spring-boot-starter-data-jpa} : persistance avec Hibernate et
    PostgreSQL (stratégie JOINED).
    \item \texttt{spring-boot-starter-security} : sécurisation des routes.
    \item \texttt{spring-boot-starter-validation} : validation des requêtes.
    \item \texttt{jjwt 0.12.6} : génération et validation des tokens JWT.
    \item \texttt{Lombok} : réduction du code boilerplate.
\end{itemize}

\subsubsection{PostgreSQL -- Base de données}

\begin{figure}[H]
    \centering
    \includegraphics[width=5cm]{PostgreSQL-Logo.png}
    \caption{PostgreSQL}
    \label{fig:postgresql}
\end{figure}

\textbf{PostgreSQL 18} est la base de données relationnelle choisie. Le schéma est
généré automatiquement par Hibernate (\texttt{ddl-auto: update}) à partir des entités
JPA annotées. Les 9 tables générées sont : \texttt{utilisateurs}, \texttt{patients},
\texttt{medecins}, \texttt{secretaires}, \texttt{dossiers\_medicaux}, \texttt{consultations},
\texttt{rendez\_vous}, \texttt{documents} et \texttt{predictions\_ia}.

\subsubsection{FastAPI -- Microservice IA}

\begin{figure}[H]
    \centering
    \includegraphics[width=5cm]{logo-fastapi.png}
    \caption{FastAPI}
    \label{fig:fastapi}
\end{figure}

Le microservice IA est développé en \textbf{Python 3.12} avec \textbf{FastAPI 0.115} et
\textbf{Uvicorn}. Il expose les endpoints suivants :

\begin{itemize}
    \item \texttt{GET /health} : vérification de l'état du service.
    \item \texttt{GET /symptoms} : liste des 132 symptômes reconnus.
    \item \texttt{POST /predict} : prédiction de maladie à partir d'une liste de
    symptômes.
    \item \texttt{POST /chat} : réponse de l'assistant médical Gemini.
\end{itemize}

Le modèle de prédiction est un \textbf{RandomForestClassifier} (scikit-learn) entraîné
sur 4\,920 échantillons synthétiques couvrant 41 maladies (accuracy : 99,9\,\%).

\subsubsection{React -- Frontend}

\begin{figure}[H]
    \centering
    \includegraphics[width=5cm]{logo_react.png}
    \caption{React}
    \label{fig:react}
\end{figure}

L'interface utilisateur est développée en \textbf{React 18} avec le bundler \textbf{Vite 6}.
Les bibliothèques principales sont :

\begin{itemize}
    \item \textbf{react-router-dom} : routage client-side avec routes protégées par
    rôle.
    \item \textbf{axios} : client HTTP avec intercepteurs JWT automatiques.
    \item \textbf{lucide-react} : bibliothèque d'icônes SVG.
\end{itemize}

% ---
\subsection{Outils}

\begin{table}[H]
\centering
\caption{Outils de développement utilisés}
\label{tab:outils}
\begin{tabular}{lll}
\toprule
\textbf{Outil} & \textbf{Version} & \textbf{Usage} \\
\midrule
Visual Studio Code & 1.100.2 & Éditeur principal \\
pgAdmin 4          & 4.x     & Administration PostgreSQL \\
Maven              & 3.9.9   & Gestionnaire de build Java \\
Git                & 2.x     & Contrôle de version \\
Postman            & 11.x    & Test des API REST \\
Overleaf           & --      & Rédaction du rapport (LaTeX) \\
\bottomrule
\end{tabular}
\end{table}

% ---
\subsection{Méthodologie}

Le développement a suivi une approche \textbf{itérative} organisée en 5 phases :

\begin{enumerate}
    \item \textbf{Phase 1 :} mise en place du backend Spring Boot -- authentification
    JWT, entités JPA, génération du schéma PostgreSQL.

    \item \textbf{Phase 2 :} développement des fonctionnalités métier -- gestion des
    utilisateurs, des rendez-vous et des dossiers médicaux.

    \item \textbf{Phase 3 :} intégration du microservice IA -- entraînement du
    modèle RandomForest, exposition via FastAPI.

    \item \textbf{Phase 4 :} développement du frontend React -- interfaces des trois
    rôles, connexion au backend via Axios.

    \item \textbf{Phase 5 :} intégration du chatbot Gemini et tests fonctionnels
    complets.
\end{enumerate}

% ---
\section{Implémentation des fonctionnalités}

\subsection{Page de connexion}

La page de connexion constitue le point d'entrée de l'application. L'utilisateur saisit
son adresse email et son mot de passe. Le backend valide les identifiants via Spring
Security, génère un token JWT valide 24~heures et retourne les informations du profil
(rôle, nom, prénom). Le frontend stocke ce token dans le \texttt{localStorage} et redirige
automatiquement l'utilisateur vers son tableau de bord selon son rôle.

\begin{figure}[H]
    \centering
    \includegraphics[width=0.7\textwidth]{screen_login.png}
    \caption{Page de connexion -- Clinique Union}
    \label{fig:login}
\end{figure}

% ---
\subsection{Interface Médecin}

Une fois connecté, le médecin accède à un espace regroupant quatre fonctionnalités
principales.

\subsubsection{Gestion des consultations}

Le médecin accède à la liste de ses consultations passées et peut en créer de nouvelles.
Lors de la création, il renseigne le motif, le diagnostic, les notes d'ordonnance et
sélectionne les symptômes observés via un champ d'autocomplétion (132 symptômes
disponibles). Dès la validation, le backend appelle automatiquement le microservice IA.

\begin{figure}[H]
    \centering
    \includegraphics[width=\textwidth]{screen_medecin_consultation.png}
    \caption{Interface Médecin -- Création de consultation avec sélection de symptômes}
    \label{fig:medecin_consultation}
\end{figure}

\subsubsection{Aide au diagnostic par IA}

Une fois les symptômes soumis, le résultat de la prédiction s'affiche sur la fiche de
consultation : maladie prédite, score de confiance en pourcentage et top-3 des pathologies
probables.

\begin{figure}[H]
    \centering
    \includegraphics[width=0.85\textwidth]{screen_medecin_ia.png}
    \caption{Interface Médecin -- Résultat de la prédiction IA}
    \label{fig:medecin_ia}
\end{figure}

\subsubsection{Gestion des documents}

Le médecin peut téléverser des documents PDF (ordonnances, résultats d'analyses)
directement depuis la fiche de consultation. Les fichiers sont stockés sur le serveur et
téléchargeables à tout moment depuis le dossier médical du patient.

\subsubsection{Planning des rendez-vous}

Le médecin consulte son planning de rendez-vous avec le statut de chaque créneau
(PLANIFIÉ, CONFIRMÉ, ANNULÉ) et les informations du patient concerné.

% ---
\subsection{Interface Patient}

L'espace patient comporte quatre sections principales.

\subsubsection{Mon dossier médical}

Le patient accède à l'historique complet de ses consultations : date, motif, diagnostic,
médecin traitant, symptômes, prédiction IA associée et documents téléchargeables.

\begin{figure}[H]
    \centering
    \includegraphics[width=\textwidth]{screen_patient_dossier.png}
    \caption{Interface Patient -- Dossier médical}
    \label{fig:patient_dossier}
\end{figure}

\subsubsection{Prise de rendez-vous}

Le patient peut demander un rendez-vous en sélectionnant un médecin, une date et un
motif. La secrétaire confirme ou refuse ensuite la demande depuis son interface.

\subsubsection{Mes rendez-vous}

Liste des rendez-vous passés et à venir avec leur statut et les informations du médecin
concerné.

\subsubsection{Assistant IA santé}

Un chatbot conversationnel propulsé par \textbf{Gemini} permet au patient de poser des
questions médicales en langage naturel. L'assistant répond en français de manière claire
et bienveillante, rappelant systématiquement qu'il ne remplace pas une consultation
médicale.

\begin{figure}[H]
    \centering
    \includegraphics[width=0.8\textwidth]{screen_patient_chatbot.png}
    \caption{Interface Patient -- Assistant IA santé (propulsé par Gemini)}
    \label{fig:patient_chatbot}
\end{figure}

% ---
\subsection{Interface Secrétaire}

La secrétaire dispose d'un tableau de bord centralisé pour la gestion administrative de
la clinique.

\begin{figure}[H]
    \centering
    \includegraphics[width=\textwidth]{screen_secretaire.png}
    \caption{Interface Secrétaire -- Tableau de bord}
    \label{fig:secretaire}
\end{figure}

\begin{itemize}
    \item \textbf{Gestion des rendez-vous :} visualisation de tous les rendez-vous,
    validation ou refus des demandes patients, modification du statut.

    \item \textbf{Gestion des patients :} liste des patients enregistrés avec accès
    à leurs informations de contact.

    \item \textbf{Suivi de l'activité :} statistiques sur les rendez-vous du jour
    et de la semaine.
\end{itemize}

\section*{Conclusion}
\addcontentsline{toc}{section}{Conclusion}

Ce chapitre a détaillé les choix techniques et l'implémentation des fonctionnalités du
système. L'architecture microservices adoptée garantit une séparation claire des
responsabilités entre le backend métier (Spring Boot), le service d'intelligence
artificielle (FastAPI) et l'interface utilisateur (React). Chaque composant peut évoluer
indépendamment, ce qui confère au système une bonne maintenabilité et une capacité
d'extension future.

% ===========================================================
% CONCLUSION GÉNÉRALE
% ===========================================================
\chapter*{Conclusion Générale}
\addcontentsline{toc}{chapter}{Conclusion Générale}
\thispagestyle{fancy}

Ce projet de fin d'année a abouti à la conception et au développement d'un système
web complet de gestion des dossiers médicaux, intégrant un module d'intelligence
artificielle pour l'aide au diagnostic. L'application offre aux médecins, patients et
secrétaires des espaces dédiés et adaptés à leurs besoins respectifs, remplaçant des
processus manuels par une solution numérique centralisée, sécurisée et accessible.

Sur le plan technique, l'architecture microservices adoptée -- combinant Spring Boot,
FastAPI et React -- s'est révélée particulièrement adaptée à ce type de système
multi-couches. L'intégration du modèle RandomForestClassifier (accuracy 99,9\,\% sur
41 maladies) et de l'assistant conversationnel Gemini constitue la fonctionnalité
différenciante du projet.

Sur le plan pédagogique, ce travail nous a permis de consolider des compétences variées :
conception orientée objet avec UML, développement backend sécurisé avec Spring Security
et JWT, apprentissage automatique avec scikit-learn, développement frontend avec React
et Vite, et gestion d'une base de données relationnelle avec PostgreSQL.

Pour l'avenir, plusieurs améliorations peuvent être envisagées :
\begin{itemize}
    \item Déploiement cloud avec Docker et Kubernetes pour une mise en production
    réelle.
    \item Intégration de notifications (email/SMS) pour les rappels de rendez-vous.
    \item Amélioration du modèle IA avec un dataset réel annoté par des médecins.
    \item Signature électronique des ordonnances.
    \item Développement d'une application mobile (React Native).
\end{itemize}

En résumé, cette solution constitue une base solide et évolutive, démontrant la
faisabilité d'un système de santé numérique abordable et intelligent pour les structures
cliniques de petite et moyenne taille.

% ===========================================================
% BIBLIOGRAPHIE ET WEBOGRAPHIE
% ===========================================================
\chapter*{Bibliographie et Webographie}
\addcontentsline{toc}{chapter}{Bibliographie et Webographie}
\thispagestyle{fancy}

\section*{1. Cours et ouvrages}

\begin{description}
    \item[NASSAR, 2025] Mahmoud NASSAR, \textit{Base de données : Modèle des données},
    ENSIAS, Année académique 2024-2025.

    \item[ETTAZI, 2025] Widad ETTAZI, \textit{Développement web}, ENSIAS,
    Année académique 2024-2025.

    \item[Walls, 2022] Craig Walls, \textit{Spring in Action}, 6\ieme{} édition,
    Manning Publications, 2022.
\end{description}

\section*{2. Logiciels utilisés}

\begin{description}
    \item[VS Code, 2024] Visual Studio Code [Logiciel], version 1.100.2.
    Disponible sur : \url{https://code.visualstudio.com}

    \item[pgAdmin, 2024] pgAdmin Development Team, \textit{pgAdmin~4} [Logiciel].
    Disponible sur : \url{https://www.pgadmin.org}

    \item[Maven, 2024] Apache Software Foundation, \textit{Apache Maven 3.9.9}.
    Disponible sur : \url{https://maven.apache.org}
\end{description}

\section*{3. Webographie}

\begin{description}
    \item[Spring Docs, 2024] Spring Framework Team,
    \textit{Spring Boot Reference Documentation}, version 3.4.x.
    Disponible sur : \url{https://docs.spring.io/spring-boot/docs/}

    \item[React Docs, 2024] Meta Platforms Inc.,
    \textit{React -- La bibliothèque pour les interfaces web}.
    Disponible sur : \url{https://react.dev}

    \item[FastAPI Docs, 2024] Sebastián Ramírez,
    \textit{FastAPI Documentation}.
    Disponible sur : \url{https://fastapi.tiangolo.com}

    \item[Kaggle Dataset, 2020] Pranay Patil,
    \textit{Disease Symptom Prediction}, Kaggle.
    Disponible sur :
    \url{https://www.kaggle.com/datasets/itachi9604/disease-symptom-description-dataset}

    \item[JWT.io, 2024] Auth0, \textit{Introduction to JSON Web Tokens}.
    Disponible sur : \url{https://jwt.io}

    \item[scikit-learn, 2024] Pedregosa et al., \textit{Scikit-learn: Machine Learning
    in Python}, JMLR 12, pp.~2825--2830, 2011.
    Disponible sur : \url{https://scikit-learn.org}

    \item[OpenRouter, 2024] OpenRouter API Documentation.
    Disponible sur : \url{https://openrouter.ai/docs}
\end{description}

\end{document}
